% engines: lualatex
% expect-warnings: 0
% expect-errors: 0
\documentclass{beamer}
\usepackage[kernel = lua]{../../latex/pgf-go}
\setbeamertemplate{navigation symbols}{}
\begin{document}

% game.sgf: 9x9; B E5, W G7, B pass, W C3, B G5.

% Slides 1-3: one diagram per slide (5 moves, step = 2).
\begin{frame}{per frame}
	\sgfdiagrams[step = 2, goban = {scale = 0.8}, layout = {per frame}]{game.sgf}
\end{frame}

% Slide 4: the text alone; slides 5-6: two diagrams per slide in a grid, after the pause.
\begin{frame}{per frame = 2, columns = 2, after a pause}
	A text before the diagrams.
	\pause
	\sgfdiagrams[step = 1, to = 4, goban = {scale = 0.6}, layout = {columns = 2, per frame = 2}]{game.sgf}
\end{frame}

% Slides 7-9: a problem, the statement alone, then three variations and then the fourth.
\begin{frame}{sgf problem}
	\sgfproblem[statement = alone, solutions = false, layout = {per frame}]{problem.sgf}
	\sgfproblem[statement = false, goban = {scale = 0.6}, layout = {columns = 3, per frame = 3}]{problem.sgf}
\end{frame}

\end{document}
